\section{Test \& Validation}
\subsection{Normal Cases}

For the three normal cases (Case 1--3), we use the simulation results obtained with \textbf{4000000.0~$\mu s$} as the main reference in the following tables. To further assess whether this simulation horizon is sufficient, we also repeated the simulations with \textbf{8000000.0~$\mu s$} and compared the maximum observed end-to-end latencies.

The additional experiment shows that \textbf{Case~2 and Case~3 are largely stable} when the simulation duration is increased from 4000000.0~$\mu s$ to 8000000.0~$\mu s$, since only very small changes are observed in a few streams. In contrast, \textbf{Case~1 still shows noticeable variation in several CBS simulation results}, which indicates that the observed maximum latency has not fully stabilized yet. Therefore, the 4000000.0~$\mu s$ results are used as the main comparison point, but they should still be interpreted as \textbf{observed maxima within the chosen simulation horizon}, rather than exact worst-case values.

\begin{table}[H]
    \centering
    \caption{Traffic Characteristics Comparison}
    \label{tab:traffic_characteristics}
    \begin{tabular}{lccc}
        \toprule
        \textbf{Metric} & \textbf{Case 1} & \textbf{Case 2} & \textbf{Case 3} \\
        \midrule
        \textbf{Topology} & ES0-SW1-ES1 & ES0-SW1-SW0-ES1 & ES0-SW1-ES1 \\
        \textbf{Total AVB Load (Mbps)} & 45.40 & 41.14 & 59.33 \\
        \textbf{Max Class A Packet Size (Bytes)} & 1012 & 1311 & 936 \\
        \textbf{Max BE Packet Size (Bytes)} & 841 & 1331 & 960 \\
        \bottomrule
    \end{tabular}
\end{table}

\subsubsection{Results}
\begin{table}[H]
    \centering
    \caption{End-to-End Latency Comparison for Case 1 (Unit: $\mu s$, duration = 4000000.0~$\mu s$)}
    \label{tab:case1}
    \begin{tabular}{cccccc}
        \toprule
        \textbf{Stream ID} & \textbf{PCP} & \textbf{Sim(CBS)} & \textbf{Sim(SP)} & \textbf{Ana(CBS)} & \textbf{Ana(SP)} \\
        \midrule
        0 & 2 & 1161.28 & 336.80 & 1493.28 & 1000.56 \\
        1 & 2 & 1373.44 & 336.80 & 1493.28 & 1000.56 \\
        2 & 2 & 1226.32 & 336.80 & 1537.68 & 1000.56 \\
        3 & 2 & 1146.48 & 336.80 & 1537.68 & 1000.56 \\
        \midrule
        4 & 1 & 630.80 & 454.16 & 1880.32 & 1629.44 \\
        5 & 1 & 1704.40 & 454.16 & 1880.32 & 1629.44 \\
        6 & 1 & 1506.08 & 454.16 & 1765.60 & 1629.44 \\
        7 & 1 & 1234.80 & 454.16 & 1765.60 & 1629.44 \\
        \midrule
        8 & 0 & 654.32 & 454.16 & NaN & 1764.00 \\
        9 & 0 & 554.32 & 454.16 & NaN & 1764.00 \\
        \bottomrule
    \end{tabular}
\end{table}

\begin{table}[H]
    \centering
    \caption{End-to-End Latency Comparison for Case 2 (Unit: $\mu s$, duration = 4000000.0~$\mu s$)}
    \label{tab:case2}
    \begin{tabular}{cccccc}
        \toprule
        \textbf{Stream ID} & \textbf{PCP} & \textbf{Sim(CBS)} & \textbf{Sim(SP)} & \textbf{Ana(CBS)} & \textbf{Ana(SP)} \\
        \midrule
        0 & 2 & 1244.41 & 451.73 & 1845.04 & 1250.13 \\
        1 & 2 & 1246.98 & 451.73 & 1845.04 & 1250.13 \\
        2 & 2 & 1240.37 & 451.73 & 1936.06 & 1250.13 \\
        3 & 2 & 1217.62 & 451.73 & 1936.06 & 1250.13 \\
        \midrule
        4 & 1 & 1319.54 & 603.49 & 2308.32 & 2024.10 \\
        5 & 1 & 1316.97 & 603.49 & 2308.32 & 2024.10 \\
        6 & 1 & 1358.45 & 603.49 & 2287.74 & 2024.10 \\
        7 & 1 & 1358.45 & 603.49 & 2287.74 & 2024.10 \\
        \midrule
        8 & 0 & 758.21 & 603.49 & NaN & 2237.06 \\
        9 & 0 & 758.18 & 603.49 & NaN & 2237.06 \\
        \bottomrule
    \end{tabular}
\end{table}

\begin{table}[H]
    \centering
    \caption{End-to-End Latency Comparison for Case 3 (Unit: $\mu s$, duration = 4000000.0~$\mu s$)}
    \label{tab:case3}
    \begin{tabular}{cccccc}
        \toprule
        \textbf{Stream ID} & \textbf{PCP} & \textbf{Sim(CBS)} & \textbf{Sim(SP)} & \textbf{Ana(CBS)} & \textbf{Ana(SP)} \\
        \midrule
        0 & 2 & 605.86 & 326.26 & 1447.36 & 950.24 \\
        1 & 2 & 605.86 & 326.26 & 1447.36 & 950.24 \\
        2 & 2 & 555.46 & 326.26 & 1406.80 & 950.24 \\
        3 & 2 & 555.46 & 326.26 & 1406.80 & 950.24 \\
        \midrule
        4 & 1 & 1177.14 & 473.78 & 2020.40 & 1713.60 \\
        5 & 1 & 1171.06 & 473.78 & 2020.40 & 1713.60 \\
        6 & 1 & 653.14 & 473.78 & 2077.28 & 1713.60 \\
        7 & 1 & 1171.06 & 473.78 & 2077.28 & 1713.60 \\
        \midrule
        8 & 0 & 610.02 & 473.78 & NaN & 1867.20 \\
        9 & 0 & 619.38 & 473.78 & NaN & 1867.20 \\
        \bottomrule
    \end{tabular}
\end{table}

\begin{itemize}
\item \textbf{Case 1:}
Case~1 has a medium AVB load and relatively small BE packets. Under SP, all AVB streams experience low latency because higher-priority queues are always served immediately. Under CBS, the AVB delays increase noticeably because transmissions are regulated by credits. BE traffic also shows moderate latency under CBS, while no analytical CBS bound is reported for BE traffic.

\item \textbf{Case 2:}
Case~2 has the lowest AVB load among the three normal cases, but it also contains the largest packets. The larger packet sizes increase blocking times and make lower-priority traffic wait longer once a transmission has started. This effect is especially visible under CBS, where both Class~B and BE streams show relatively large simulated latencies despite the lower total AVB load.

\item \textbf{Case 3:}
Case~3 has the highest AVB load, but the packets are smaller than in Case~2. As a result, blocking intervals are shorter, and the BE delays remain lower than in Case~2. However, the CBS delays of some Class~B streams are still large, which shows that total load and packet size both influence the observed latency.
\end{itemize}

\subsubsection{Evaluation \& Discussion}

Based on the results of the three normal test cases, the following observations can be made regarding the behavior of Credit-Based Shaper (CBS) and Strict Priority (SP):

\begin{itemize}
\item \textbf{Safety Verification for AVB Traffic:}
For the AVB streams in the normal cases, the simulated maximum latencies remain below the corresponding analytical bounds in both SP and CBS. This indicates that the analytical model provides a safe upper bound for AVB traffic in these scenarios.

\item \textbf{SP Gives Lower Latency for AVB Traffic:}
Strict Priority (SP) consistently yields lower latency for Class~A and Class~B streams than CBS. This is expected because SP always serves the highest-priority non-empty queue immediately, while CBS intentionally delays AVB transmissions through credit regulation.

\item \textbf{CBS Does Not Provide a BE Analytical Bound:}
In the CBS analytical results, the entries for Best Effort traffic are reported as \texttt{NaN}. This means that no analytical CBS bound is provided for BE traffic, because the lecture CBS analysis used in this project applies only to AVB traffic.

\item \textbf{Impact of Traffic Pattern, Not Only Total Load:}
The delay behavior is influenced not only by the total offered AVB load, but also by packet sizes and traffic composition. Case~2 demonstrates this clearly: although it has the lowest AVB load, it still produces relatively large simulated delays because the larger packets create longer blocking intervals.

\item \textbf{Duration Comparison for Normal Cases:}
To evaluate whether the simulation window was long enough, we compared the results obtained with 4000000.0~$\mu s$ and 8000000.0~$\mu s$. Case~2 and Case~3 showed only very small differences, which suggests that the 4000000.0~$\mu s$ horizon already provides a good approximation for these two cases. However, Case~1 still exhibited noticeable changes in several simulated CBS maxima, indicating that the observed worst-case phasing has not fully stabilized yet. Therefore, the 4000000.0~$\mu s$ results are useful for comparison, but they should still be interpreted cautiously as simulation-based observations rather than exact worst-case delays.
\end{itemize}

\subsection{Starvation Cases}

Test Case~4 is designed as an \textbf{extreme overload scenario} to evaluate the traffic isolation and protection capabilities of Credit-Based Shaper (CBS) versus Strict Priority (SP).

\begin{itemize}
    \item \textbf{Topology}: ES0 $\rightarrow$ SW0 $\rightarrow$ ES1 (single bottleneck link).
    \item \textbf{Link Capacity}: 100 Mbps.
    \item \textbf{Traffic Load}:
\item \textbf{Class A (PCP 2)}: 3 streams $\times$ 40 Mbps = \textbf{120 Mbps}.
\item \textbf{Class B (PCP 1)}: 2 streams $\times$ 8 Mbps = \textbf{16 Mbps}.
\item \textbf{Best Effort (PCP 0)}: 2 streams $\times$ 8 Mbps = \textbf{16 Mbps}.
\item \textbf{Total Load}: $120 + 16 + 16 = \mathbf{152 \text{ Mbps}}$, which exceeds the 100 Mbps link capacity.
\end{itemize}

\subsubsection{Results}

\begin{table}[H]
    \centering
    \caption{End-to-End Latency Comparison for Case 4 (Overload Scenario, Unit: $\mu s$, duration = 4000000.0~$\mu s$)}
    \label{tab:case4}
    \begin{tabular}{cccccc}
        \toprule
        \textbf{Stream ID} & \textbf{PCP} & \textbf{Sim(CBS)} & \textbf{Sim(SP)} & \textbf{Ana(CBS)} & \textbf{Ana(SP)} \\
        \midrule
        0 & 2 & 2327214.06 & 661034.06 & 2040.00 & 1320.00 \\
        1 & 2 & 2327414.06 & 660974.06 & 2040.00 & 1320.00 \\
        2 & 2 & 2327394.06 & 660914.06 & 2040.00 & 1320.00 \\
        \midrule
        3 & 1 & 814.06 & 0.00 & 1560.00 & $\infty$ \\
        4 & 1 & 814.06 & 0.00 & 1560.00 & $\infty$ \\
        \midrule
        5 & 0 & 1014.06 & 0.00 & NaN & $\infty$ \\
        6 & 0 & 1014.06 & 0.00 & NaN & $\infty$ \\
        \bottomrule
    \end{tabular}
\end{table}

\subsubsection{Evaluation \& Discussion}

\begin{itemize}

\item \textbf{The Overload Case Leads to Partial Divergence:}
In Case~4, the offered traffic exceeds the link capacity, resulting in persistent backlog growth. However, the divergence is not uniform across all traffic classes. The high-priority Class~A streams exhibit continuously increasing delays, while lower-priority traffic behaves differently depending on the scheduling mechanism.

\item \textbf{SP Causes Complete Starvation of Lower-Priority Traffic:}
Under SP scheduling, Class~A traffic fully occupies the link due to its high load. As a result, Class~B and BE streams receive no service during the simulation, leading to observed latencies of 0.00~$\mu s$. This indicates \textbf{complete starvation}, where lower-priority queues are never scheduled.

\item \textbf{CBS Provides Isolation by Sacrificing Class A Latency:}
Under CBS, credit regulation prevents Class~A traffic from continuously occupying the link. This allows Class~B and BE traffic to be transmitted with small and stable delays (around 800--1000~$\mu s$). However, this protection comes at the cost of \textbf{unbounded delay growth for Class~A streams}, whose latency increases significantly as backlog accumulates.

\item \textbf{CBS Analytical Results for BE Remain Undefined:}
As in the normal cases, the CBS analytical results for BE are reported as \texttt{NaN}, since the analytical model used here only applies to AVB traffic.

\item \textbf{Duration Comparison Confirms Divergence Behavior:}
When increasing the simulation duration from 4000000.0~$\mu s$ to 8000000.0~$\mu s$, the latency of Class~A streams under both CBS and SP approximately doubles, confirming persistent backlog growth. In contrast, the delays of Class~B and BE streams under CBS remain stable, while under SP they remain unscheduled. This demonstrates that CBS enforces traffic isolation, whereas SP leads to starvation of lower-priority traffic.

\end{itemize}

\section{Conclusion}

Based on the four experimental cases, the following key insights regarding the behavior of Credit-Based Shaper (CBS) can be summarized:

\begin{itemize}
\item \textbf{SP Minimizes AVB Delay but Strongly Favors High Priorities:}
Across the normal cases, SP consistently gives lower delays for Class~A and Class~B traffic than CBS. This is because SP always serves the highest-priority non-empty queue immediately.

\item \textbf{CBS Introduces Additional Delay Through Credit Regulation:}
CBS regulates AVB transmissions using credits, which increases delay compared with SP even in normal operating conditions. This effect is clearly visible in Cases~1--3.

\item \textbf{The Traffic Pattern Matters as Much as the Total Load:}
Differences in packet size and traffic composition strongly affect the observed delays. Case~2 shows that lower total AVB load does not necessarily imply lower latency when packet sizes are larger and blocking becomes more significant.

\item \textbf{Simulation Horizon Must Be Interpreted Carefully:}
Comparing 4000000.0~$\mu s$ and 8000000.0~$\mu s$ shows that the observed simulated maxima can still change with a longer horizon. Case~2 and Case~3 are largely stable, but Case~1 still exhibits noticeable variation in several CBS simulation results. Therefore, the simulation results should be interpreted as observed maxima within the selected time window, not as guaranteed exact worst-case delays.

\item \textbf{Overload Causes Divergence and Starvation Effects:}
In the overload scenario, delays of high-priority traffic grow continuously due to backlog accumulation. Under SP, lower-priority traffic can suffer from complete starvation and may receive no service at all. In contrast, CBS enforces isolation by ensuring service for lower-priority traffic, at the cost of significantly increased delays for higher-priority streams.
\end{itemize}
